\documentclass[12pt]{article}
\usepackage[utf8]{inputenc}
\usepackage[T1]{fontenc}
\usepackage[margin=1in]{geometry}
\usepackage{setspace}
\usepackage{xcolor}
\usepackage[colorlinks=true, urlcolor=blue!60!black, linkcolor=blue!60!black]{hyperref}
\setlength{\parindent}{0em}
\setlength{\parskip}{0.6em}
\renewcommand{\baselinestretch}{1.0}
\begin{document}
\begin{center}
{\LARGE \textit{Shurtleff v. City of Boston}, 596 U.S. 243 (2022)\par}
\vspace{0.5em}
{\large Case Brief --- Two-Page Edition\par}
\end{center}

\section{Memory Jogger}
After approving 284 private flag raisings at City Hall without a single refusal, Boston rejected a Christian flag; the Court held 9--0 that the program was private speech because Boston did not control the flags' messages, so the refusal was viewpoint discrimination (\href{https://supreme.justia.com/cases/federal/us/596/20-1800/}{596 U.S. 243}).

\section{Detailed Case Facts}
Boston let private groups raise flags of their choosing on the third flagpole at City Hall Plaza, approving about 50 flags in 284 ceremonies with no written policy. In 2017 Harold Shurtleff of Camp Constitution asked to raise the ``Christian flag.'' A city commissioner refused, fearing an Establishment Clause violation.

\section{Procedural History}
The District Court granted Boston summary judgment, holding the flags government speech, and the First Circuit affirmed. The Supreme Court reversed.

\section{Judicial Votes}
Reversed, 9--0. Breyer wrote for Roberts, Sotomayor, Kagan, Kavanaugh, and Barrett; Kavanaugh concurred; Alito (with Thomas and Gorsuch) and Gorsuch (with Thomas) concurred in the judgment.

\section{Holding}
Under a holistic inquiry into history, public perception, and active government control, Boston's program was private speech, and refusing a flag for its religious viewpoint violated the Free Speech Clause.

\section{Analysis of Opinions}
Breyer found that history favored Boston and perception was inconclusive, but control was decisive: Boston shaped the flags' messages ``not at all,'' showing a ``come-one-come-all attitude'' unlike \textit{Summum} or \textit{Walker}. With government speech gone, rejecting a flag because it promoted a religion was viewpoint discrimination. Kavanaugh stressed that religion may not be treated as second-class. Alito would limit government speech to messages the government purposefully communicates through its agents. Gorsuch blamed \textit{Lemon} for Boston's error.

\section{Examples: Future Applications}
\textbf{Same-Side:} a town that approves every holiday display may not reject an atheist group's sign. \textbf{Opposite-Side:} a written closed-list flag policy makes the pole government speech. \textbf{Fence-Sitter:} a juried mural program approving about 90 percent of proposals without content rules.

\section{Critique}
\textit{Shurtleff} may encourage cities to level down by closing forums, and, as Alito warned, the control factor can reward the very censorship it should deter.

\section{Key Quotations}
\begin{quote}
``In answering these questions, we conduct a holistic inquiry designed to determine whether the government intends to speak for itself or to regulate private expression. Our review is not mechanical; it is driven by a case's context rather than the rote application of rigid factors.''
\end{quote}
\noindent --- Justice Breyer, Opinion of the Court, 596 U.S. at 252.

\end{document}
